% $Id: abstract.tex 1789 2010-09-28 16:30:23Z jabriffa $

\chapter*{Abstract}

When a computation is outsourced to an untrusted party, two things must hold:
the result must be correct, and any sensitive inputs must remain private. A
zero-knowledge virtual machine (zkVM) addresses the correctness side by
generating a cryptographic proof that a fixed program ran correctly, but
confidentiality still depends on the cryptographic primitive embedded in that
program. Choosing the right primitive is therefore a concrete engineering
decision, and one that cannot be resolved by circuit-level theory alone --- the
zkVM adds its own execution overhead that can invert expected rankings.

This dissertation studies the practical cost of symmetric cryptography inside a
zero-knowledge virtual machine using RISC Zero, with the aim of identifying
which cipher and mode combination is best suited for authenticated encryption in
this setting. The work starts with operation-level benchmarking and baseline
cipher benchmarking for AES-192, Salsa20, and LowMC, all measured through a
common host/guest harness that emits structured JSON artifacts for later
aggregation.

The measurement workflow is designed for repeatability: benchmark targets are
config-driven, trials are logged per run, and aggregation emits machine-readable
summaries suitable for chapter-level evidence inserts. The final main and
operations report campaigns are commit-locked, clean-tree, interleaved
five-trial runs; older supporting snapshots are identified where they remain
single-trial.

The initial hypothesis was that LowMC should be favorable because of reduced
nonlinearity. In this implementation setting, measured behavior did not match
that expectation. The project therefore includes optimization work on LowMC
(bitsliced SWAR substitution, matrix hot-path changes, and precompute-oriented
host/guest inputs) and an AES optimization workspace for S-box backend
experimentation.

The final implemented system path is authenticated AES-192 in the zkVM:
AES-192-CTR encryption with a truncated HMAC-SHA256 tag. Security-sensitive
cryptographic work remains inside the guest so public outputs are bound to the
RISC Zero receipt.


Security analysis is presented as a game-based proof sketch rather than a full
formal reduction. The main game is an IND-CPA challenge where the adversary also
receives the HMAC tag and RISC Zero receipt; the sketch relies on nonce-respecting
AES-CTR, a PRF-style treatment of truncated HMAC tags, and proof zero-knowledge
for non-journaled witness data.

The evaluation provides artifact-backed results. In the latest main report run,
optimised LowMC remains about \(33.33\times\) slower than AES on median proving
time (484.279 s vs 14.530 s), while supporting LowMC baseline snapshots show that
local optimisation helped but did not reverse the ranking. For the final
authenticated AES pipeline, adding HMAC increases user-cycle work at every tested
block size; matched wall-clock total-time deltas are not significant after
correction. These results are therefore reported with explicit statistical and
reproducibility caveats rather than as cross-zkVM-general claims.
